% DERIVED from formal_layer.json — read-only, do not edit.
\begin{assumptionv}{ass:iid-sampling}[Independent observed records]
The observed records \(O_1,\ldots,O_n\) are independent and identically distributed under the observed margin of \(P\).
\end{assumptionv}

\begin{assumptionv}{ass:consistency}[Consistency of observed outcomes]
\(Y=Y(A)\) almost surely.
\end{assumptionv}

\begin{assumptionv}{ass:conditional-exchangeability}[Conditional exchangeability of treatment]
Under \(P\), \((Y(0),Y(1))\perp A\mid X\).
\end{assumptionv}

\begin{definitionv}{synth_1}[Observations from potential outcomes]
For \(d\in\mathbb N\), let \(P\) be a full-data causal law on \((X,A,Y(0),Y(1))\), where \(X\) ranges over a \(d\)-element alphabet and \(A,Y(0),Y(1)\in\{0,1\}\). The associated full observation is defined atomwise by
\[
(X,A,Y(0),Y(1))\longmapsto
\bigl(X,A,Y(A),Y(0),Y(1)\bigr),
\qquad
Y(A)=
\begin{cases}
Y(1),& A=1,\\
Y(0),& A=0.
\end{cases}
\]
The observed margin of \(P\) is the law of \((X,A,Y(A))\).
\end{definitionv}

\begin{definitionv}{synth_2}[Observed covariate-cell mass]
For a discrete law \(P\) on a covariate alphabet of size \(d\), the observed mass of cell \(j\) is
\[
p_j := P(X=j).
\]
\end{definitionv}

\begin{definitionv}{synth_3}[Product law and propensity]
For a discrete observed-data law \(P\) on a \(d\)-cell covariate alphabet and a sample size \(n\), the product law is the joint law of \(n\) independent records \(O_i=(X_i,A_i,Y_i)\) drawn from \(P\). For each occupied cell \(j\), write \(p_j=P(X=j)>0\) and define its treatment propensity by \(e_j:=P(A=1\mid X=j)\).
\end{definitionv}

\begin{assumptionv}{ass:overlap}[Fixed treatment overlap]
For every \(j\in[d]\) with \(p_j>0\), the treatment propensity satisfies
\[
\epsilon\le e_j\le 1-\epsilon.
\]
\end{assumptionv}

\begin{definitionv}{def:model-class}[Causal law class \(\mathcal M_{d,\epsilon}\)]
\(\mathcal M_{d,\epsilon}\), for each integer \(d\ge2\), is the class of laws \(P\) satisfying:
\begin{itemize}
\item \textbf{(Covariate alphabet.)} \(X\in[d]\).
\item \textbf{(Binary variables.)} \(A,Y(0),Y(1)\in\{0,1\}\).
\item \textbf{(Identification conditions.)} \cref{obj:ass:consistency,obj:ass:conditional-exchangeability,obj:ass:overlap}.
\end{itemize}
\end{definitionv}

\begin{definitionv}{def:budget-policy-class}[Budget policy class \(\Pi_b(P)\)]
\[
\Pi_b(P)=\left\{\pi\in[0,1]^d:\sum_{j=1}^d p_j\pi_j\le b\right\}.
\]
\end{definitionv}

\begin{definitionv}{synth_7}[Control-policy value]
For a full-data law \(P\), let \(p_j=P(X=j)\) and let \(\mu_{0j}=E_P[Y(0)\mid X=j]\) on occupied cells, with \(\mu_{0j}=0\) on null cells. The population value of assigning every cell to control is
\[
B(P)=\sum_{j=1}^d p_j\mu_{0j}.
\]
\end{definitionv}

\begin{definitionv}{synth_4}[Cell potential-outcome means]
For a full-data causal law \(P\) on a \(d\)-cell covariate alphabet, an arm \(a\in\{0,1\}\), and a cell \(j\), let \(p_j=P(X=j)\). On occupied cells, where \(p_j>0\), define
\[
\mu_{aj}=\mathbb E_P\!\left[Y(a)\mid X=j\right].
\]
On null cells, where \(p_j=0\), set \(\mu_{aj}=0\).
\end{definitionv}

\begin{definitionv}{synth_5}[Cell treatment effect]
For a full-data causal law \(P\) on a \(d\)-cell covariate alphabet and a cell \(j\), define the cell treatment effect by
\[
\tau_j := \mu_{1j}-\mu_{0j},
\]
where \(\mu_{aj}\) is the potential-outcome mean for arm \(a\) in cell \(j\).
\end{definitionv}

\begin{definitionv}{def:budget-value}[Budget value $V_b(P)$ and curve $V(P)$]
The maximum population value at budget \(b\) and the budget-indexed oracle-value curve on \(I\) are
\[
V_b(P)=B(P)+\max_{\pi\in\Pi_b(P)}\sum_{j=1}^d p_j\pi_j\tau_j,
\qquad
V(P)=\bigl(V_b(P):b\in I\bigr).
\]
\end{definitionv}

\begin{definitionv}{synth_9}[Arm and total masses of a four-vector]
For \(u=(u_{ay}:(a,y)\in\{0,1\}^2)\in\mathbb R_+^4\), define the mass in treatment arm \(a\) and the total mass by
\[
s_a(u)=u_{a0}+u_{a1},\qquad
s(u)=s_0(u)+s_1(u)=\sum_{a,y\in\{0,1\}}u_{ay}.
\]
\end{definitionv}

\begin{definitionv}{synth_6}[Potential atom space]
For a covariate alphabet of size \(d\in\mathbb N\), a potential atom is an element of
\[
\{j\in\mathbb N:j<d\}\times\{0,1\}\times\{0,1\}\times\{0,1\}.
\]
\end{definitionv}

\begin{definitionv}{def:dual-process}[Dual threshold process $F_P(\lambda)$]
For \(u\in\mathbb R_+^4\), define the stabilized arm contributions, cell contribution, and dual threshold process by
\[
g_a(u)=\frac{s(u)u_{a1}}{\max\{s_a(u),\epsilon s(u)\}},
\qquad
f_\lambda(u)=g_0(u)+\bigl[g_1(u)-g_0(u)-\lambda s(u)\bigr]_+,
\qquad
F_P(\lambda)=\sum_{j=1}^d f_\lambda(q_j).
\]
When all four cell masses vanish (\(u=0\)), the quotient defining \(g_a(u)\) takes the value zero: \(g_a(0)=0\).
\end{definitionv}

\begin{definitionv}{synth_12}[Binding half-budget capacity]
We say the half-budget constraint is binding for \(P\) when there exists an optimal policy \(\pi^\star\in\Pi_{1/2}(P)\) such that
\[
V_{1/2}(P)=B(P)+\sum_{j=1}^d p_j\pi^\star_j\tau_j
\quad\text{and}\quad
\sum_{j=1}^d p_j\pi^\star_j=\frac12.
\]
\end{definitionv}

\begin{propositionv}{prop:dual-representation}[Dual representation of budget value]
Let \(P\) be a full-data causal law with observed margin. Suppose:
\begin{itemize}
\item \textbf{Alphabet size.} The known covariate alphabet has size \(d\ge 2\).
\item \textbf{Overlap parameter.} The fixed overlap parameter satisfies \(0<\epsilon<1/2\).
\item \textbf{Model membership.} \(P\in\mathcal M_{d,\epsilon}\) as defined in \cref{obj:def:model-class}.
\end{itemize}
For \(0\le\lambda\le1\), write \(F_P(\lambda)=\sum_{j=1}^{d}f_\lambda(q_j)\) for the dual threshold process formed from the observed margin of \(P\). Then, for every \(b\in[0,1]\), the budget value \(V_b(P)\) of \cref{obj:def:budget-value} satisfies
\[
V_b(P)=\inf_{0\le\lambda\le1}\{b\lambda+F_P(\lambda)\}.
\]
Moreover, for every process \(\widehat F\) bounded on \([0,1]\),
\[
\sup_{b\in[0,1]}
\left|
\inf_{0\le\lambda\le1}\{b\lambda+\widehat F(\lambda)\}
-V_b(P)
\right|
\le
\sup_{\lambda\in[0,1]}
\left|\widehat F(\lambda)-F_P(\lambda)\right|.
\]
\end{propositionv}

\begin{definitionv}{synth_11}[Finite observed sample]
For \(n,d\in\mathbb N\), the observed sample \(O^{(n)}=(O_i)_{i=1}^n\) consists of \(n\) observed records \(O_i=(X_i,A_i,Y_i)\), each with \(X_i\) in the \(d\)-element covariate alphabet.
\end{definitionv}

\begin{definitionv}{def:curve-risk}[Curve risk $\mathfrak R_{n,d,\epsilon,b_0}^{\mathrm{curve}}$]
For \(b_0\in[0,1/2]\), set \(I=[b_0,1-b_0]\). The all-procedure expected squared supremum risk is
\[
\mathfrak R_{n,d,\epsilon,b_0}^{\mathrm{curve}}
=
\inf_{\widehat V}
\sup_{P\in\mathcal M_{d,\epsilon}}
E_P\!\left[
\sup_{b\in I}
\left|\widehat V_b(O^{(n)})-V_b(P)\right|^2
\right],
\]
where the infimum ranges over all measurable maps from the observed sample into bounded functions on \(I\).
\end{definitionv}

\begin{theoremv}{thm:matched-curve-frontier}[Matched curve minimax frontier]*
Fix an overlap parameter \(\epsilon\in(0,1/2)\). There exist constants \(0<c_\epsilon\le C_\epsilon<\infty\) such that, for every \(n\ge1\), \(d\ge2\), and \(b_0\in[0,1/2]\), the curve minimax risk of \cref{obj:def:curve-risk} satisfies
\[
c_\epsilon\min\left\{1,\frac{d}{n\log(ed)}\right\}
\le \mathfrak R_{n,d,\epsilon,b_0}^{\mathrm{curve}}
\le C_\epsilon\min\left\{1,\frac{d}{n\log(ed)}\right\}.
\]
For every measurable scalar estimator \(\widehat V_{1/2}\) based on the \(n\) observed records, there is a full-data law \(P\) with the following properties:
\begin{itemize}
\item \textbf{(Causal model.)} \(P\in\mathcal M_{d,\epsilon}\), and each occupied cell has \(p_j>0\), \(\tau_j\in[1/8,3/8]\), and \(\tau_j>0\).
\item \textbf{(Values and binding capacity.)} \(V_1(P)=1/2\), and there is a policy \(\pi\in[0,1]^d\) such that
\[
\sum_{j=1}^{d}p_j\pi_j=\frac12,
\qquad
V_{1/2}(P)=B(P)+\sum_{j=1}^{d}p_j\pi_j\tau_j.
\]
\item \textbf{(Scalar risk.)} Under \(n\) independent observations from the observed margin of \(P\),
\[
\mathbb E_P\!\left[\bigl(\widehat V_{1/2}-V_{1/2}(P)\bigr)^2\right]
\ge c_\epsilon\min\left\{1,\frac{d}{n\log(ed)}\right\}.
\]
\end{itemize}
When \(b_0=0\), every \(P\in\mathcal M_{d,\epsilon}\) also satisfies the full-budget identity
\[
V_1(P)=B(P)+\sum_{j=1}^{d}p_j\max\{\tau_j,0\},
\]
where \(V_b(P)\) is the budget value of \cref{obj:def:budget-value}.
\end{theoremv}
% CAUSALSMITH-CITED-SCOPE-BEGIN thm:matched-curve-frontier
\verificationfootnotetext{\textbf{Formalization scope.} This result uses the published conclusion \citep{JiaoHanWeissman2018} (Theorem 3 and the unknown-both-distributions reduction following Theorem 6). The cited source proof is not formalized here: Lean verifies its use and all remaining steps. If that cited conclusion is false, the portions of this result that depend on it are not certified.}
% CAUSALSMITH-CITED-SCOPE-END thm:matched-curve-frontier

\begin{theoremv}{thm:capacity-active-lower}[Capacity-active risk lower bound]*
Fix an overlap parameter \(0<\epsilon<1/2\). There exists \(c_\epsilon>0\) such that, for every \(n\geq1\), \(d\geq2\), every collection of observed-data laws satisfying \cref{obj:ass:iid-sampling} for each full-data law, and every measurable scalar estimator \(T\) based on the \(n\) observed records, there is a full-data law \(P\) satisfying:
\begin{itemize}
\item \textbf{(Model class.)} \(P\in\mathcal M_{d,\epsilon}\).
\item \textbf{(Full-budget value.)} \(V_1(P)=1/2\), with \(V_b(P)\) as in \cref{obj:def:budget-value}.
\item \textbf{(Effects in occupied cells.)} For every cell \(j\), \(p_j>0\) implies \(1/8\leq\tau_j\leq3/8\).
\item \textbf{(Binding capacity.)} The half-budget constraint is binding for \(P\).
\end{itemize}
For this law,
\[
E_P\!\left[(T-V_{1/2}(P))^2\right]
\geq c_\epsilon\min\left\{1,\frac{d}{n\log(ed)}\right\}.
\]
\end{theoremv}
% CAUSALSMITH-CITED-SCOPE-BEGIN thm:capacity-active-lower
\verificationfootnotetext{\textbf{Formalization scope.} This result uses the published conclusion \citep{JiaoHanWeissman2018} (Theorem 3 and the unknown-both-distributions reduction following Theorem 6). The cited source proof is not formalized here: Lean verifies its use and all remaining steps. If that cited conclusion is false, the portions of this result that depend on it are not certified.}
% CAUSALSMITH-CITED-SCOPE-END thm:capacity-active-lower

\begin{propositionv}{prop:fixed-alphabet-reduction}[Fixed alphabet rate and identity]*
Fix the overlap parameter \(\epsilon\) and the known covariate alphabet size \(d\) under the following conditions:
\begin{itemize}
\item \textbf{(Overlap.)} \(0<\epsilon<1/2\).
\item \textbf{(Alphabet.)} \(d\ge 2\).
\end{itemize}
For the curve minimax risk in \cref{obj:def:curve-risk}, there are constants \(c,C\) with \(0<c\le C\) such that, for every \(n\ge1\) and every \(b_0\in[0,1/2]\),
\[
\frac{c}{n}\le
\mathfrak R_{n,d,\epsilon,b_0}^{\mathrm{curve}}
\le\frac{C}{n}.
\]

Moreover, let \(P\) belong to the causal model class \(\mathcal M_{d,\epsilon}\) in \cref{obj:def:model-class}. Write \(p_j\) for the mass of cell \(j\), \(\mu_{aj}\) for its arm-\(a\) potential-outcome mean, and \(\tau_j=\mu_{1j}-\mu_{0j}\) for its treatment effect. If some \(c>0\) satisfies \(\tau_j=c\) in every cell with \(p_j>0\), then, for every \(b_0\in[0,1/2]\) and every \(b\in[b_0,1-b_0]\), the budget value in \cref{obj:def:budget-value} satisfies
\[
V_b(P)=B(P)+cb,
\qquad
B(P)=\sum_{j=1}^{d}p_j\mu_{0j}.
\]
\end{propositionv}
% CAUSALSMITH-CITED-SCOPE-BEGIN prop:fixed-alphabet-reduction
\verificationfootnotetext{\textbf{Formalization scope.} This result uses the published conclusion \citep{JiaoHanWeissman2018} (Theorem 3 and the unknown-both-distributions reduction following Theorem 6). The cited source proof is not formalized here: Lean verifies its use and all remaining steps. If that cited conclusion is false, the portions of this result that depend on it are not certified.}
% CAUSALSMITH-CITED-SCOPE-END prop:fixed-alphabet-reduction

\begin{definitionv}{synth_8}[Observed category-cell masses]
For a covariate cell \(j\) and a treatment-outcome category \(\zeta=(a,y)\in\{0,1\}^2\), define
\[
q_{\zeta j}=P(X=j,A=a,Y=y).
\]
These four masses are the components of the observed cell vector \(q_j=(q_{\zeta j}:\zeta\in\{0,1\}^2)\).
\end{definitionv}

\begin{algorithmv}{def:jf-estimator}[Jackson–factorial estimator $\widehat V_b^{\mathrm{JF}}$]
Given \(O^{(n)}\), \(d\), and \(I\), the estimator \(\widehat V_b^{\mathrm{JF}}\) uses the tuning quantities \(H_0,\kappa,D_0,L,m,K\) defined below.
\begin{enumerate}
\item Set
\[
H_0=4096,\qquad
\kappa=\frac1{512},\qquad
D_0=16,\qquad
L=\log(ed),\qquad
m=\frac n8,\qquad
K=\max\left\{2,\left\lfloor\frac L{512}\right\rfloor\right\},
\qquad
\mathcal Z=\{0,1\}^2.
\]
Index treatment-outcome categories by \(\zeta=(a,y)\in\mathcal Z\).

\item If \(d<D_0\), form the empirical category-cell masses and threshold process:
\[
\widehat q_{\zeta j}
=
\frac1n\sum_{i=1}^n
\mathbf1\{X_i=j,\ (A_i,Y_i)=\zeta\},
\qquad
\widehat q_j=(\widehat q_{\zeta j})_{\zeta\in\mathcal Z},
\qquad
\widehat F(\lambda)=\sum_{j=1}^d f_\lambda(\widehat q_j).
\]
Proceed to the final step.

\item If \(d\geq D_0\) and \(n<d/L\), output
\[
\widehat V_b^{\mathrm{JF}}=\frac12,\qquad b\in I.
\]
In the remaining branch, \(d\geq D_0\) and \(n\geq d/L\), draw \(M\sim\operatorname{Pois}(n/4)\) independently of \(O^{(n)}\). On \(\{M>n\}\), set
\[
\widetilde F^{\mathrm{cap}}(\lambda)=0,\qquad \lambda\in[0,1].
\]
On \(\{M\leq n\}\), uniformly permute the records, retain the first \(M\), and independently mark each retained record pilot or evaluation with probability \(1/2\) each. Define the pilot and evaluation counts by
\[
N'_{\zeta j}
=
\sum_{i=1}^{M}
\mathbf1\{X_i^*=j,\ (A_i^*,Y_i^*)=\zeta,\ i\text{ is marked pilot}\},
\qquad
N_{\zeta j}
=
\sum_{i=1}^{M}
\mathbf1\{X_i^*=j,\ (A_i^*,Y_i^*)=\zeta,\ i\text{ is marked evaluation}\},
\]
where \(O_i^*=(X_i^*,A_i^*,Y_i^*)\) denotes the \(i\)th record in the uniform permutation.

\item On \(\{M\leq n\}\), for every \(j\) and \(\zeta\), set
\[
\delta=\frac Lm,\qquad
c_{\zeta j}=\frac{N'_{\zeta j}}m,\qquad
h_{\zeta j}=H_0\left\{\sqrt{c_{\zeta j}\delta}+\delta\right\},
\]
\[
\ell_{\zeta j}=\max\{0,c_{\zeta j}-h_{\zeta j}\},
\qquad
u_{\zeta j}=c_{\zeta j}+h_{\zeta j},
\qquad
c^0_{\zeta j}=\frac{\ell_{\zeta j}+u_{\zeta j}}2,
\qquad
R_{\zeta j}=\frac{u_{\zeta j}-\ell_{\zeta j}}2,
\qquad
S_j=\sum_{\zeta\in\mathcal Z}R_{\zeta j}.
\]
All radii \(R_{\zeta j}\) are strictly positive. Let \(c_j^0=(c^0_{\zeta j})_{\zeta\in\mathcal Z}\) and \(R_j=(R_{\zeta j})_{\zeta\in\mathcal Z}\). Define the normalized positive order-four Jackson kernel by
\[
J_K(t)=C_K\left\{\frac{\sin(Kt/2)}{\sin(t/2)}\right\}^{4},
\qquad -\pi\leq t\leq\pi,
\qquad
\int_{-\pi}^{\pi}J_K(t)\,dt=1,
\]
where the quotient takes its continuous value at zero and \(C_K>0\) is the unique normalizing constant. Set \(D=2(K-1)\). For each \(j\) and \(\lambda\), define the unique tensor polynomial \(P_{j,\lambda}\), of coordinatewise degree at most \(D\), through
\[
P_{j,\lambda}(c_j^0+R_j\odot\cos\vartheta)
=
\int_{[-\pi,\pi]^4}
f_\lambda\!\left(c_j^0+R_j\odot\cos(\vartheta+t)\right)
\prod_{\zeta\in\mathcal Z}J_K(t_\zeta)\,dt,
\]
and define its coefficients \(\beta_{j,\alpha}(\lambda)\) by the unique expansion
\[
P_{j,\lambda}(u)-f_\lambda(c_j^0)
=
\sum_{\alpha\in\{0,\ldots,D\}^4}
\beta_{j,\alpha}(\lambda)
\prod_{\zeta\in\mathcal Z}
\left(\frac{u_\zeta-c^0_{\zeta j}}{R_{\zeta j}}\right)^{\alpha_\zeta}.
\]
For integers \(h\geq0\), define
\[
U_h(N;z)
=
\sum_{t=0}^h {h\choose t}(-z)^{h-t}\frac{(N)_t}{m^t},
\qquad
(N)_t=N(N-1)\cdots(N-t+1),
\qquad
(N)_0=1.
\]
The cell estimate, its clipped version, and the auxiliary threshold process are
\[
Z_j(\lambda)
=
f_\lambda(c_j^0)+
\sum_{\alpha\in\{0,\ldots,D\}^4}
\beta_{j,\alpha}(\lambda)
\prod_{\zeta\in\mathcal Z}
\frac{U_{\alpha_\zeta}(N_{\zeta j};c^0_{\zeta j})}
{R_{\zeta j}^{\alpha_\zeta}},
\]
\[
T_j(\lambda)
=
f_\lambda(c_j^0)+
\Pi_{[-d^{1/4}S_j,d^{1/4}S_j]}
\!\left(Z_j(\lambda)-f_\lambda(c_j^0)\right),
\qquad
\widetilde F^{\mathrm{cap}}(\lambda)
=
\sum_{j=1}^d T_j(\lambda).
\]
Here
\[
\Pi_{[r,s]}(x)=\min\{s,\max\{r,x\}\}.
\]
In this branch, average the auxiliary process conditional on the observed sample:
\[
\widehat F(\lambda)
=
E\!\left\{\widetilde F^{\mathrm{cap}}(\lambda)\mid O^{(n)}\right\}.
\]

\item In the empirical and Jackson–factorial branches, report
\[
\widehat V_b^{\mathrm{JF}}
=
\Pi_{[0,1]}
\left(
\min_{0\leq\lambda\leq1}
\{b\lambda+\widehat F(\lambda)\}
\right),
\qquad b\in I.
\]
The output uses the minimum value and requires no optimizer selection. If an implementation stores an optimizer, it stores the least element of the nonempty compact argmin set. These branch rules, the zero process on \(\{M>n\}\), and the least-minimizer convention apply under every data or treatment-effect tie.
\end{enumerate}
\end{algorithmv}

\begin{theoremv}{thm:curve-upper}[Uniform budget-curve upper bound]
Fix the overlap parameter \(\epsilon\). Suppose:
\begin{itemize}
\item \textbf{(Overlap.)} \(0<\epsilon<1/2\).
\item \textbf{(Sample and alphabet.)} \(n\geq 1\) and \(d\geq 2\).
\item \textbf{(Budget interval.)} \(b_0\in[0,1/2]\), and \(I=[b_0,1-b_0]\).
\item \textbf{(Sampling.)} For each causal law \(P\in\mathcal M_{d,\epsilon}\), the \(n\) observed records follow \cref{obj:ass:iid-sampling} for the observed margin of \(P\).
\end{itemize}
There is a constant \(C_\epsilon>0\), depending only on \(\epsilon\), such that the estimator in \cref{obj:def:jf-estimator} and the budget value in \cref{obj:def:budget-value} satisfy
\[
\sup_{P\in\mathcal M_{d,\epsilon}}
E_P\!\left[
\sup_{b\in I}
\left|\widehat V_b^{\mathrm{JF}}-V_b(P)\right|^2
\right]
\leq
C_\epsilon\min\left\{1,\frac{d}{n\log(ed)}\right\}.
\]
\end{theoremv}

\begin{lemmav}{lem:bv-lipschitz}[Bounded variation Lipschitz bound]
For \(0<\epsilon<1/2\), there is a finite constant \(C_\epsilon\ge 0\) such that, for every \(u,v\in\mathbb R_+^4\),
\[
\|f_{\cdot}(u)-f_{\cdot}(v)\|_{BV([0,1])}
\le C_\epsilon\|u-v\|_1,
\]
where \(\|h\|_{BV([0,1])}=|h(0)|+\operatorname{TV}_{[0,1]}(h)\).
\end{lemmav}

\begin{definitionv}{def:process-handle}[Uniform process handle $\mathsf H_\epsilon$]
For fixed \(0<\epsilon<1/2\), the uniform process handle \(\mathsf H_\epsilon\) is the set of constants \(C>0\) satisfying the five groups of bounds below for the single pilot-local Jackson polynomial and clipping rule of \(\widehat V^{\mathrm{JF}}\). Membership requires the bounds simultaneously for every \(n,d\) in the Jackson–factorial branch, every observed categorical law \(P\) on the known \(d\)-cell alphabet, every pilot realization, every cell \(j\), and every \(\lambda\in[0,1]\).

Set
\[
L=\log(ed),\qquad
m=\frac n8,\qquad
K=\max\left\{2,\left\lfloor\frac L{512}\right\rfloor\right\},
\qquad
S_j=\sum_{\zeta\in\mathcal Z}R_{\zeta j},
\qquad
v_j=\sqrt{\frac{p_jL}{m}}+\frac Lm.
\]
For the expectations and \(L_2\) norms below, use the uncapped marked-Poisson experiment in which all pilot and evaluation counts \(N'_{\zeta j}\) and \(N_{\zeta j}\) are mutually independent with distribution \(\operatorname{Pois}(m q_{\zeta j})\). Define
\[
G_j
=
\bigcap_{\zeta\in\mathcal Z}
\left\{
|c_{\zeta j}-q_{\zeta j}|
\leq \frac{h_{\zeta j}}4
\right\},
\qquad
\|h\|_{BV([0,1])}
=
|h(0)|+\operatorname{TV}_{[0,1]}(h).
\]
Evaluate \(T_j\), \(\beta_{j,\alpha}\), and \(R_{\zeta j}\) in this experiment. The membership bounds are:
\begin{itemize}
\item \textbf{Coefficient envelope and exponential bound.}
\[
\sum_\alpha\|\beta_{j,\alpha}\|_{BV([0,1])}
\leq C S_j e^{12K},
\qquad
S_j^2 e^{24K+16K^2/L}
\leq C d^{1/16}S_j^2.
\]

\item \textbf{Centered process on the pilot event.}
\[
E\left\|
\sum_{j=1}^d
\left[
\mathbf1_{G_j}\{T_j-f_\cdot(q_j)\}
-
E\mathbf1_{G_j}\{T_j-f_\cdot(q_j)\}
\right]
\right\|_\infty^2
\leq C\frac d{mL}.
\]

\item \textbf{Cellwise bias.}
\[
\sup_{0\leq\lambda\leq1}
\left|ET_j(\lambda)-f_\lambda(q_j)\right|
\leq
C\left\{\sqrt{\frac{p_j}{mL}}+\frac1{mL}\right\}.
\]

\item \textbf{Cellwise envelope outside the pilot event.}
\[
\left\|
\mathbf1_{G_j^c}
\sup_{0\leq\lambda\leq1}
|T_j(\lambda)-f_\lambda(q_j)|
\right\|_{L_2}
\leq
C d^{1/4}e^{-10L}v_j.
\]

\item \textbf{Centered process outside the pilot event.}
\[
E\left\|
\sum_{j=1}^d
\left[
\mathbf1_{G_j^c}\{T_j-f_\cdot(q_j)\}
-
E\mathbf1_{G_j^c}\{T_j-f_\cdot(q_j)\}
\right]
\right\|_\infty^2
\leq C\frac d{mL}.
\]
\end{itemize}
\end{definitionv}

\begin{lemmav}{lem:pilot-radius-second-moment}[Second moment of pilot radii]
Let \(n\geq 1\), \(d\geq 16\), and \(m=n/8\). Let \(P\) be a categorical law on the \(d\)-cell alphabet, with category-cell masses \(q_{\zeta j}\) and cell masses \(p_j=\sum_{\zeta\in\mathcal Z}q_{\zeta j}\), and fix a cell \(j\). In the ideal Poisson experiment, the pilot counts \(N'_{\zeta j}\) are independent with means \(m q_{\zeta j}\). Let \(S_j=\sum_{\zeta\in\mathcal Z}R_{\zeta j}\) be the sum of the four pilot radii in \cref{obj:def:jf-estimator}, and set
\[
L=\log(ed),\qquad H_0=4096,\qquad
v_j=\sqrt{\frac{p_jL}{m}}+\frac{L}{m}.
\]
Then \(S_j^2\) is integrable under the ideal Poisson law, and
\[
\mathbb E[S_j^2]\leq 32H_0^2v_j^2.
\]
\end{lemmav}

\begin{lemmav}{lem:jackson-boundary-adaptive}[Boundary-adaptive Jackson approximation]
Let \(P_\lambda\) be the local polynomial for \(f_\lambda\) of \cref{obj:def:dual-process}, formed by the normalized order-four Jackson construction in \cref{obj:def:jf-estimator} with pilot center and radii replaced by \(c^0\) and \(R\), and with the value \(f_\lambda(c^0)\) added to the normalized polynomial evaluation. Suppose that:
\begin{itemize}
\item \textbf{(Parameters.)} The overlap parameter satisfies \(\epsilon>0\), the shadow price satisfies \(\lambda\in[0,1]\), and \(K\) is a positive integer.
\item \textbf{(Rectangle.)} For every \(\zeta\in\mathcal Z\), \(R_\zeta>0\) and \(c^0_\zeta-R_\zeta\geq0\).
\item \textbf{(Evaluation point.)} For every \(\zeta\in\mathcal Z\), \(|y_\zeta-c^0_\zeta|\leq R_\zeta\).
\end{itemize}
Then
\[
\bigl|P_\lambda(y)-f_\lambda(y)\bigr|
\leq
32\bigl\{3(1+\epsilon^{-1})+1\bigr\}
\sum_{\zeta\in\mathcal Z}
\left\{
\frac{\sqrt{(y_\zeta-c^0_\zeta+R_\zeta)(c^0_\zeta+R_\zeta-y_\zeta)}}{K}
+\frac{R_\zeta}{K^2}
\right\}.
\]
\end{lemmav}

\begin{lemmav}{lem:pilot-bad-score-moment}[Bad pilot score moment]
Let \(n,d\geq 1\), let \(P\) be an observed categorical law on the \(d\)-cell alphabet, and fix a cell \(j\). Write \(q_{\zeta j}\) for its four category-cell masses and \(p_j=\sum_{\zeta\in\mathcal Z}q_{\zeta j}\). Set \(m=n/8\) and \(L=\log(ed)\).

Suppose the four pilot counts are independent, with \(N'_{\zeta j}\sim\operatorname{Pois}(mq_{\zeta j})\) for each \(\zeta\in\mathcal Z\). Form \(c_{\zeta j}\) and \(h_{\zeta j}\) as in \cref{obj:def:jf-estimator}, and let \(G_j\) be the quarter-halfwidth pilot event of \cref{obj:def:process-handle}. Define
\[
B_j=\sum_{\zeta\in\mathcal Z}\bigl(|c_{\zeta j}-q_{\zeta j}|+h_{\zeta j}\bigr),
\qquad
v_j=\sqrt{\frac{p_jL}{m}}+\frac{L}{m}.
\]
For each \(t\in\{1,2,4\}\),
\[
\mathbb E\!\left[\mathbf 1_{G_j^c}B_j^{\,t}\right]
\leq
4^t\,10^{80(t+1)}e^{-20L}v_j^{\,t}.
\]
\end{lemmav}

\begin{lemmav}{lem:bv-maximal-cell-paths}[Maximal bound for centered cell paths]
Let \(n,d\) be nonnegative integers, let \(\epsilon\in\mathbb R\), and let \(P\) be a discrete law on the \(d\)-cell alphabet with category-cell masses \(q_{\zeta j}\). In the ideal count experiment of \cref{obj:def:process-handle}, the pilot and evaluation counts \(N'_{\zeta j}\) and \(N_{\zeta j}\) are mutually independent with \(\operatorname{Pois}(mq_{\zeta j})\) laws, where \(m=n/8\). Choose the pilot events uniformly across cells: either \(E_j=G_j\) for every \(j\), or \(E_j=G_j^c\) for every \(j\), with \(G_j\) as in \cref{obj:def:process-handle}. For each cell \(j\), define the path on \([0,1]\) by
\[
W_j(\lambda)=\mathbf 1_{E_j}\bigl\{T_j(\lambda)-f_\lambda(q_j)\bigr\},
\]
where \(T_j\) is the clipped cell estimate of \cref{obj:def:jf-estimator} and \(f_\lambda\) is defined in \cref{obj:def:dual-process}. Assume:

\begin{itemize}
\item \textbf{(Regularity.)} For every count table and cell \(j\), \(W_j\) is continuous in \(\lambda\). For each \(j\), the count-table map into \(C([0,1])\) is measurable and Bochner integrable.
\item \textbf{(Finite variation.)} For each \(j\), \(\operatorname{TV}_{[0,1]}(W_j)<\infty\) almost surely.
\item \textbf{(Second moments.)} For each \(j\), \(\bigl(\|W_j\|_\infty+\operatorname{TV}_{[0,1]}(W_j)\bigr)^2\) is integrable.
\item \textbf{(Independence.)} The paths \(W_1,\ldots,W_d\) are independent.
\end{itemize}

Then
\[
\mathbb E\left\|\sum_{j=1}^{d}\bigl(W_j-\mathbb EW_j\bigr)\right\|_\infty^2
\leq 16384\sum_{j=1}^{d}
\mathbb E\left[\bigl(\|W_j\|_\infty+\operatorname{TV}_{[0,1]}(W_j)\bigr)^2\right].
\]
\end{lemmav}

\begin{lemmav}{lem:capped-marked-count-identity}[Capped marked-count integral identity]
Fix \(n,d\in\mathbb N\) and an observed categorical law \(P\) on the \(d\)-cell alphabet, with category-cell masses \(q_{\zeta j}\). Let \(O^{(n)}=(O_1,\ldots,O_n)\sim P^{\otimes n}\). For \(M\in\{0,\ldots,n\}\), a permutation \(\sigma\) of the \(n\) record indices, and pilot/evaluation marks \(\omega\in\{\mathrm{pilot},\mathrm{evaluation}\}^n\), write \((N',N)(O^{(n)},\sigma,M,\omega)\) for the pilot and evaluation count tables formed from the first \(M\) permuted records according to \cref{obj:def:jf-estimator}. Define
\[
w_n(M)=\frac{\Pr\{\operatorname{Pois}(n/4)=M\}}{n!\,2^n}.
\]
In the uncapped marked-Poisson experiment, \(M\sim\operatorname{Pois}(n/4)\); conditional on \(M\), draw \(M\) independent records from \(P\) and mark each independently as pilot or evaluation with probability \(1/2\), yielding count tables \((N',N)\).

For every real-valued function \(g\) of pairs of count tables that is integrable under the ideal law of \cref{obj:def:process-handle}, whose pilot and evaluation category-cell counts are independent \(\operatorname{Pois}((n/8)q_{\zeta j})\) variables,
\[
\mathbb E_{O^{(n)}\sim P^{\otimes n}}\left[
\sum_{M=0}^{n}\sum_{\sigma}\sum_{\omega\in\{\mathrm{pilot},\mathrm{evaluation}\}^n}
w_n(M)\,
g\bigl((N',N)(O^{(n)},\sigma,M,\omega)\bigr)
\right]
=
\mathbb E_{\mathrm{uncapped}}\left[
\mathbf 1\{M\leq n\}\,g(N',N)
\right].
\]
\end{lemmav}

\begin{theoremv}{thm:integrated-process-certificate}[Uniform integrated process bounds]
Fix an overlap parameter \(0<\epsilon<1/2\). There is a single constant \(C_\epsilon>0\), belonging to the uniform process handle in \cref{obj:def:process-handle}, for which the following conclusions hold whenever:

\begin{itemize}
\item \textbf{(Sample and alphabet.)} \(n\geq1\), \(d\geq16\), and \(d/\log(ed)\leq n\).
\item \textbf{(Sampling.)} The law \(\mu\) of the \(n\) observed records satisfies \cref{obj:ass:iid-sampling} for the observed margin of \(P\).
\item \textbf{(Causal law.)} The full-data law \(P\) belongs to \(\mathcal M_{d,\epsilon}\) of \cref{obj:def:model-class}.
\end{itemize}

Write \(L=\log(ed)\), \(m=n/8\), and \(v_j=\sqrt{p_jL/m}+L/m\). Simultaneously for every cell \(j\) and every \(\lambda\in[0,1]\), the absolute bias of the uncapped Jackson cell estimate of \(f_\lambda(q_j)\) is at most
\[
C_\epsilon\left\{\sqrt{\frac{p_j}{nL}}+\frac{1}{nL}\right\}.
\]
The expected squared supremum of the centered good-pilot error sum is at most \(C_\epsilon d/(nL)\). The uncentered bad-pilot cell envelope satisfies
\[
\left\|\mathbf 1_{G_j^c}\sup_{\lambda\in[0,1]}
\left|T_j(\lambda)-f_\lambda(q_j)\right|\right\|_{L_2}
\leq C_\epsilon d^{1/4}e^{-10L}v_j,
\]
and the expected squared supremum of the centered bad-pilot error sum is at most \(C_\epsilon d/(mL)\). Finally, the ideal and averaged versions of the dual threshold process in \cref{obj:def:jf-estimator} each have expected squared supremum error relative to \(F_P\) at most \(C_\epsilon d/(nL)\). The same \(C_\epsilon\) applies to all these bounds.
\end{theoremv}

\begin{lemmav}{lem:empirical-threshold-process-bounded}[Bounded empirical threshold process]
For any \(n,d\in\mathbb N\), any real overlap parameter \(\epsilon\), and any sample \(O^{(n)}\) of \(n\) observations on the \(d\)-cell alphabet, let
\[
\widehat F(\lambda)=\sum_{j=1}^{d}f_\lambda(\widehat q_j)
\]
be the empirical threshold process of \cref{obj:def:jf-estimator}, with \(f_\lambda\) as in \cref{obj:def:dual-process}. There exists \(B\in\mathbb R\) such that
\[
|\widehat F(\lambda)|\leq B
\qquad\text{for every }\lambda\in[0,1].
\]
\end{lemmav}

\begin{lemmav}{lem:averaged-threshold-process-bounded}[Bounded averaged threshold process]
Fix any observed sample \(O^{(n)}\) of \(n\) records from a \(d\)-cell alphabet. Suppose:
\begin{itemize}
\item \textbf{(Overlap parameter.)} \(\epsilon>0\).
\item \textbf{(Sample size.)} \(n\geq 1\).
\item \textbf{(Alphabet size.)} \(d\geq 1\).
\end{itemize}
The conditionally averaged threshold process \(\widehat F(\lambda)\) of \cref{obj:def:jf-estimator}, whose Poisson\((n/4)\) weights are summed over \(M=0,\ldots,n\), is bounded on \([0,1]\): there exists \(B\in\mathbb R\) such that
\[
|\widehat F(\lambda)|\leq B
\qquad\text{for every }\lambda\in[0,1].
\]
\end{lemmav}

\begin{definitionv}{def:paired-family}[Paired causal family $\mathcal P_k^{\mathrm{pair}}$]
The paired causal family \(\mathcal P_k^{\mathrm{pair}}\) consists of causal completions constructed as follows. For \(k\geq1\), choose \(r=(r_1,\ldots,r_k)\) in the simplex and \(\theta_i\in[-1/8,1/8]\). For each \(i\), create cells \((i,+)\) and \((i,-)\) with masses \(r_i/2\), and set
\[
e_{i,\sigma}=\frac12,\qquad
\mu_{0,i,\sigma}=\frac14,\qquad
\mu_{1,i,\sigma}=\frac12+\sigma\theta_i,
\qquad
\sigma\in\{+1,-1\}.
\]
Generate the potential outcomes independently of \(A\) conditional on \(X\).
\end{definitionv}

\begin{definitionv}{synth_10}[Probability simplex on \(k\) categories]
For an integer \(k\geq1\), define
\[
\Delta_k=\left\{r=(r_1,\ldots,r_k)\in[0,1]^k:\sum_{i=1}^k r_i=1\right\}.
\]
\end{definitionv}

\begin{propositionv}{prop:paired-capacity-identity}[Paired capacity and sample identity]
Let \(k\geq 1\) and \(0<\epsilon<1/2\). The following statements hold.
\begin{itemize}
\item \textbf{(Capacity identity.)} For every \(P\in\mathcal P^{\mathrm{pair}}_k\) and every pair-mass vector \(r\in\Delta_k\) and contrast vector \(\theta\in[-1/8,1/8]^k\) for which \(P\) satisfies the paired-completion conditions of \cref{obj:def:paired-family}, the law \(P\) belongs to \(\mathcal M_{2k,\epsilon}\) of \cref{obj:def:model-class}. Every occupied cell has treatment effect \(\tau_j\in[1/8,3/8]\). The half-population budget is attained by a policy \(\pi\in\Pi_{1/2}(P)\) satisfying
\[
\sum_j p_j\pi_j=\frac12,
\qquad
V_{1/2}(P)=B(P)+\sum_j p_j\pi_j\tau_j.
\]
For the budget values of \cref{obj:def:budget-value},
\[
V_{1/2}(P)
=\frac38+\frac12\sum_{i=1}^k r_i|\theta_i|
=\frac18+F_P\!\left(\frac14\right),
\qquad
V_1(P)=\frac12.
\]

\item \textbf{(Exact sample reduction.)} For every integer \(n\geq0\), there is a Markov kernel from pairs of \(n\)-record samples on the \(k\)-point alphabet to \(n\) observed records on the \(2k\)-cell alphabet with the following property. For every pair \(R,S\in\Delta_k\), the simplex of \(k\)-point laws, set
\[
r_i=\frac{R_i+S_i}{2},
\qquad
\theta_i=
\begin{cases}
0, & R_i+S_i=0,\\[2pt]
\dfrac{R_i-S_i}{8(R_i+S_i)}, & R_i+S_i>0.
\end{cases}
\]
The paired law determined by \((r,\theta)\) belongs to \(\mathcal P^{\mathrm{pair}}_k\). Applying the kernel to independent \(n\)-record samples from \(R\) and \(S\) yields exactly the law of \(n\) independent observed records from the observed margin of that paired law.
\end{itemize}
\end{propositionv}

\begin{lemmav}{lem:two-sample-l1-lower}[Two-sample \(L_1\) lower bound]*
There is a universal constant \(c>0\) such that, for every integer \(k\geq 2\) and \(1\leq n\leq k^2\), the minimax squared-error risk for estimating \(\|R-S\|_1\) from \(n\) independent observations from each of \(R,S\in\Delta_k\), the probability simplex on \(k\) categories, satisfies
\[
\inf_{\widehat L}\sup_{R,S\in\Delta_k}
E_{R,S}\!\left[(\widehat L-\|R-S\|_1)^2\right]
\geq c\min\left\{1,\frac{k}{n\log(ek)}\right\}.
\]
\end{lemmav}
% CAUSALSMITH-CITED-SCOPE-BEGIN lem:two-sample-l1-lower
\verificationfootnotetext{\textbf{Formalization scope.} This result uses the published conclusion \citep{JiaoHanWeissman2018} (Theorem 3 and the unknown-both-distributions reduction following Theorem 6). The cited source proof is not formalized here: Lean verifies its use and all remaining steps. If that cited conclusion is false, the portions of this result that depend on it are not certified.}
% CAUSALSMITH-CITED-SCOPE-END lem:two-sample-l1-lower
